\documentclass[12pt]{article}
\usepackage[T1]{fontenc}
\usepackage[utf8]{inputenc}
\usepackage{lmodern}
\usepackage{textcomp}
\usepackage{enumitem}
\usepackage{geometry}
\geometry{margin=1in}
\begin{document}
\begin{center}
Answer Key - Constitutional Law - Graduate Level Assessment 17 \
\small (Answers are ordered exactly as in the question paper.)
\end{center}

\section*{Multiple Choice}
\begin{enumerate}[label=\arabic*.]
\item \textbf{Answer:} A
\\ \textit{Options:} A) not hearsay.; B) hearsay, but admissible as an admission.; C) hearsay, but admissible as a dying declaration.; D) hearsay not within any recognized exception.
\\ \textit{Note:} The testimony is not hearsay because it is offered for the truth of the matter asserted, which is that the witness hired someone to commit the murder, thereby providing a potentially exculpatory statement relevant to the defendant's case.
\item \textbf{Answer:} D
\\ \textit{Options:} A) Nominal; B) Punitive; C) Rescission; D) Compensatory
\\ \textit{Note:} Compensatory damages are designed to reimburse the nonbreaching party for losses suffered due to the breach, thereby restoring them to the position they would have been in had the contract been fully performed.
\item \textbf{Answer:} D
\\ \textit{Options:} A) Pay just compensation for the property; B) Terminate the regulation and pay damages that occurred while regulation was in effect; C) Pay double damages if not addressed in 30 days.; D) a or b
\\ \textit{Note:} The correct answer is "a or b" because, under the Fifth Amendment, if a regulation constitutes a taking of private property, the government is required to either provide just compensation to the property owner or demonstrate a legitimate public purpose for the taking.
\item \textbf{Answer:} A
\\ \textit{Options:} A) Gender; B) Interstate travel; C) Privacy; D) Alienage
\\ \textit{Note:} Gender classifications are subject to an intermediate scrutiny standard of review rather than strict scrutiny, which is reserved for classifications based on race or fundamental rights.
\item \textbf{Answer:} B
\\ \textit{Options:} A) upon court order.; B) upon request.; C) upon request and showing of good cause.; D) under no circumstances.
\\ \textit{Note:} The correct answer is "upon request" because constitutional law mandates that defendants have the right to request certain information, including a list of prospective jurors, to ensure a fair trial and the ability to challenge potential jurors effectively.
\end{enumerate}

\section*{True / False}
\begin{enumerate}[label=\arabic*.]
\setcounter{enumi}{5}
\item \textbf{Answer:} True
\\ \textit{Options:} A) True; B) False
\\ \textit{Note:} The Sixth Amendment right to counsel applies to criminal prosecutions and formal adversarial proceedings, and since investigative surveillance does not constitute such a proceeding, an indigent person does not have the right to counsel at that stage.
\item \textbf{Answer:} True
\\ \textit{Options:} A) True; B) False
\\ \textit{Note:} Contracts for crops and timber to be severed are classified under the UCC because they are considered goods under the definition provided in UCC Article 2, which governs the sale of tangible personal property.
\item \textbf{Answer:} False
\\ \textit{Options:} A) True; B) False
\\ \textit{Note:} The Rule in Shelley's Case operates to merge a life estate and a remainder into a fee simple when the remainder is granted to the heirs of the life tenant, thereby preventing a remainder in the grantor's heirs and allowing the property to pass directly to the life tenant's heirs instead of reverting to the grantor.
\item \textbf{Answer:} True
\\ \textit{Options:} A) True; B) False
\\ \textit{Note:} Common law burglary is defined as the unlawful entry into a building with the intent to commit a crime inside, regardless of whether the entry occurs during the day or night, thus making the statement true.
\item \textbf{Answer:} True
\\ \textit{Options:} A) True; B) False
\\ \textit{Note:} The necessary and proper clause, while granting Congress the ability to enact laws that carry out its enumerated powers, cannot independently justify federal legislation without the existence of an explicit enumerated power in the Constitution to support that law.
\end{enumerate}

\section*{Long Form Response}
\begin{enumerate}[label=\arabic*.]
\setcounter{enumi}{10}
\item \textbf{Answer:} In constitutional law, the warrant requirement is subject to several exceptions, including exigent circumstances, consent, search incident to arrest, and the automobile exception. However, the concept of cold pursuit, which refers to the immediate apprehension of a suspect without a warrant based on the belief that they are committing or have committed a crime, is not recognized as a valid exception. This distinction is critical because it underscores the necessity of balancing law enforcement needs with Fourth Amendment protections against unreasonable searches and seizures. Cold pursuit raises concerns about the potential for abuse and the erosion of constitutional safeguards, emphasizing the importance of judicial oversight in maintaining the integrity of the warrant requirement. Thus, while various exceptions allow for warrantless searches, cold pursuit remains outside this framework, reinforcing the necessity of adhering to constitutional protections.
\\ \textit{Note:} correct because it accurately identifies specific exceptions to the warrant requirement while emphasizing that cold pursuit lacks legal recognition as a valid exception, thereby highlighting the importance of constitutional protections against unreasonable searches and seizures.
\item \textbf{Answer:} The legal status of a jogger who enters a convenience store to use the bathroom during a run is that of a public invitee. Property owners have a duty to ensure the safety of invitees on their premises, which includes providing reasonable access to facilities such as restrooms. This status implies that the jogger has a right to enter the store for a legitimate purpose, and the store owner cannot unjustly deny access or create unsafe conditions. The implications within the context of property rights and public access highlight the balance between private property rights and the societal expectation of providing basic amenities to the public, particularly in commercial spaces. Thus, the jogger's status as a public invitee reinforces the idea that businesses have an obligation to accommodate individuals who may require access to their facilities.
\\ \textit{Note:} correct because it identifies the jogger as a public invitee, which grants them the right to access the store's facilities and imposes a legal obligation on the property owner to ensure their safety, reflecting the balance between property rights and public access.
\end{enumerate}

\vfill
\noindent\textit{End of Answer Key}
\end{document}